\documentclass[conference]{IEEEtran}
\IEEEoverridecommandlockouts
\usepackage[utf8]{inputenc}
\usepackage[T1]{fontenc}
\usepackage{amsmath,amssymb,amsfonts}
\usepackage{algorithmic}
\usepackage{graphicx}
\usepackage{textcomp}
\usepackage{xcolor}
\usepackage{booktabs}
\usepackage{longtable}
\usepackage{array}
\usepackage{hyperref}
\usepackage{url}
\usepackage{enumitem}
\usepackage{microtype}
\usepackage{caption}
\usepackage{float}
\usepackage{setspace}

\hypersetup{
    colorlinks=true,
    linkcolor=blue,
    citecolor=blue,
    urlcolor=blue
}

\makeatletter
\def\BibTeX{{\rm B\kern-.05em{\sc i\kern-.025em b}\kern-.08em
    T\kern-.1667em\lower.7ex\hbox{E}\kern-.125emX}}
\makeatother

\begin{document}

\title{Digitalisasi Layanan Melalui Aplikasi SIP-HAPUS: Solusi Terpadu
Untuk Optimalisasi Proses Penghapusan Regident Kendaraan Bermotor
Di Jawa Tengah\\
{\footnotesize \textsuperscript{*}Note: Sub-titles are not captured for
https://ieeexplore.ieee.org and should not be used}
}

\author{
\IEEEauthorblockN{1\textsuperscript{st} Bagaskoro Saputro}
\IEEEauthorblockA{\textit{Computer Science, BINUS @Semarang Campus}\\
\textit{BINUS University}\\
Semarang, Indonesia\\
bagaskoro.saputro@binus.ac.id}
\and
\IEEEauthorblockN{2\textsuperscript{nd} Galih Dea Pratama}
\IEEEauthorblockA{\textit{Computer Science, BINUS University}\\
\textit{BINUS University}\\
Semarang, Indonesia\\
}
\and
\IEEEauthorblockN{3\textsuperscript{rd} Dimas Elang Setyoko}
\IEEEauthorblockA{\textit{Computer Science, BINUS University}\\
\textit{BINUS University}\\
Semarang, Indonesia\\
}
\and
\IEEEauthorblockN{4\textsuperscript{th} Canggih Gelar Setyo Adhi}
\IEEEauthorblockA{\textit{Computer Science, BINUS University}\\
\textit{BINUS University}\\
Semarang, Indonesia\\
}
\and
\IEEEauthorblockN{5\textsuperscript{th} Dian Martha Nurrul Amanah}
\IEEEauthorblockA{\textit{Digital Business, BINUS University}\\
\textit{BINUS University}\\
Semarang, Indonesia\\
}
\and
\IEEEauthorblockN{6\textsuperscript{th} Mutiara Adinda}
\IEEEauthorblockA{\textit{Digital Business, BINUS University}\\
\textit{BINUS University}\\
Semarang, Indonesia\\
}
\and
\IEEEauthorblockN{7\textsuperscript{th} Freysia Chandra Saliman}
\IEEEauthorblockA{\textit{Digital Business, BINUS University}\\
\textit{BINUS University}\\
Semarang, Indonesia\\
}
\and
\IEEEauthorblockN{8\textsuperscript{th} Early Achiril Putra}
\IEEEauthorblockA{\textit{Computer Science, BINUS University}\\
\textit{BINUS University}\\
Semarang, Indonesia\\
}
\and
\IEEEauthorblockN{9\textsuperscript{th} Lintang Nathaniela Pribadi}
\IEEEauthorblockA{\textit{Computer Science, BINUS University}\\
\textit{BINUS University}\\
Semarang, Indonesia\\
}
\and
\IEEEauthorblockN{10\textsuperscript{th} Arfa Naufal Azizan}
\IEEEauthorblockA{\textit{Computer Science, BINUS University}\\
\textit{BINUS University}\\
Semarang, Indonesia\\
}
\and
\IEEEauthorblockN{11\textsuperscript{th} Maximilian Otto Mukti Aji}
\IEEEauthorblockA{\textit{Digital Business, BINUS University}\\
\textit{BINUS University}\\
Semarang, Indonesia\\
}
}

\maketitle

\begin{abstract}
\textbf{Latar Belakang} --- Proses penghapusan registrasi dan
identifikasi kendaraan bermotor di Provinsi Jawa Tengah saat ini
melibatkan banyak instansi pemerintah, yaitu Samsat, Kepolisian
Daerah (Polda), Badan Pendapatan Daerah (Bapenda), dan PT Jasa Raharja.
Koordinasi lintas instansi yang masih bergantung pada dokumen fisik
dan alur verifikasi manual menimbulkan berbagai permasalahan struktural,
di antaranya proses administrasi yang panjang dan tidak terintegrasi,
ketidaktransparanan dalam monitoring status pengajuan, potensi
miskomunikasi antar instansi, serta belum adanya dokumentasi aktivitas
yang terpusat. Kondisi ini sejalan dengan konsep \emph{governance gap}
dalam tata kelola layanan publik berbasis digital (Janssen et al., 2020),
di mana ketidakmampuan sistem dalam menyediakan transparansi dan
keterlacakan proses secara signifikan menurunkan kualitas layanan.

\textbf{Tujuan} --- Penelitian ini bertujuan untuk merancang dan
mengembangkan sebuah platform monitoring terintegrasi berbasis web yang
mampu mendigitalisasi seluruh alur pengajuan penghapusan regident
kendaraan bermotor, menyediakan fitur \emph{tracking} status secara
\emph{real-time}, dashboard monitoring, manajemen Surat Keputusan (SK),
notifikasi otomatis, serta \emph{role-based access control} (RBAC)
untuk menjamin keamanan dan relevansi akses data.

\textbf{Metode} --- Penelitian menggunakan kerangka \emph{Design Science
Research} (DSR) sebagaimana dikembangkan oleh Peffers et al.~(2007),
yang terdiri dari enam tahapan utama, yaitu identifikasi masalah
(\emph{problem identification}), penetapan tujuan solusi (\emph{define
objectives of a solution}), perancangan dan pengembangan artefak
(\emph{design and development}), demonstrasi (\emph{demonstration}),
evaluasi (\emph{evaluation}), dan komunikasi (\emph{communication}).
Pengembangan sistem menggunakan framework Laravel dengan pola arsitektur
\emph{Model-View-Controller} (MVC), basis data MySQL, \emph{containerization}
dengan Docker, serta \emph{deployment} pada infrastruktur \emph{cloud computing}
Google Cloud Platform (GCP) mencakup Compute Engine, Cloud Storage, dan
Cloud SQL. Evaluasi dilakukan melalui \emph{expert heuristic evaluation}
berdasarkan 10 prinsip usability Nielsen (Nielsen, 1994).

\textbf{Hasil} --- Platform berhasil dikembangkan dengan tujuh modul
utama yang saling terintegrasi: Dashboard Monitoring, Tracking Pengajuan,
Manajemen Pengajuan, Manajemen Surat Keputusan, Activity Log, Notifikasi
Status, dan Role-Based Access Control. \emph{Heuristic evaluation}
terhadap tujuh halaman utama aplikasi menunjukkan bahwa 82,9\% dari
seluruh item evaluasi mendapat skor ``Good'', 12,9\% ``Minor'', dan
hanya 1,4\% ``Major'' yang terkonsentrasi pada aspek dokumentasi dalam
aplikasi (H10). Simulasi \emph{System Usability Scale} (SUS) menghasilkan
skor 71,5 yang termasuk dalam Grade B dengan kategori ``Good'' dan
melampaui rata-rata industri sistem \emph{e-government} sebesar 68.
Estimasi \emph{task completion rate} mencapai 92,3\% dengan rata-rata
waktu penyelesaian 2,2 menit per tugas dari 13 skenario yang diuji.

\textbf{Kesimpulan} --- Platform SIP-HAPUS berhasil mengintegrasikan
proses bisnis pencabutan pajak kendaraan bermotor ke dalam sistem digital
yang transparan, \emph{real-time}, dan aman. Seluruh luaran yang ditargetkan
dalam program Hibah Internal PkM Universitas Bina Nusantara, yaitu publikasi
Scopus, poster ilmiah, \emph{enrichment} pengajaran, dan teknologi/inovasi,
berada dalam tahap pengerjaan sesuai dengan \emph{roadmap} yang telah
ditetapkan.

\textbf{Kata Kunci}: Penghapusan Regident, Design Science Research,
Laravel, Heuristic Evaluation, Sistem Monitoring, E-Government
\end{abstract}

\begin{IEEEkeywords}
Penghapusan Regident, Design Science Research, Laravel, Heuristic Evaluation,
Sistem Monitoring, E-Government
\end{IEEEkeywords}

\section{Introduction}
\label{sec:introduction}

\subsection{Latar Belakang Masalah}
\label{subsec:latar-belakang}

Proses pencabutan pajak kendaraan bermotor di Jawa Tengah merupakan sebuah
rangkaian administrasi yang melibatkan kompleksitas koordinasi lintas instansi
pemerintah. Alur pengajuan penghapusan registrasi dan identifikasi kendaraan
bermotor dimulai dari wajib pajak yang mengajukan permohonan, kemudian melewati
proses verifikasi di Samsat, verifikasi lanjutan di Kepolisian Daerah (Polda),
persetujuan dari Badan Pendapatan Daerah (Bapenda) dan PT Jasa Raharja, hingga
penerbitan Surat Keputusan (SK) oleh masing-masing instansi terkait. Setiap
tahapan dalam alur tersebut saat ini masih dilakukan secara manual dan
terpisah-pisah antar instansi, sehingga berpotensi memperlambat keseluruhan
proses administrasi.

Berdasarkan analisis terhadap alur bisnis yang berjalan, terdapat beberapa
permasalahan mendasar yang perlu mendapatkan penanganan secara sistematis.
Pertama, proses administrasi yang panjang dan tidak terintegrasi menjadi kendala
utama karena setiap instansi memiliki sistem dan prosedur kerja yang berdiri
sendiri tanpa adanya keterhubungan data secara langsung. Kedua,
ketidaktransparanan dalam monitoring status pengajuan menyebabkan wajib pajak
tidak memiliki akses secara \emph{real-time} untuk memantau sejauh mana proses
pengajuan mereka telah berjalan. Ketiga, potensi miskomunikasi dan keterlambatan
antar instansi semakin tinggi karena data yang belum terpusat masih mengandalkan
dokumen fisik dan koordinasi informal yang rentan terhadap kesalahan dan duplikasi
informasi. Keempat, belum adanya dokumentasi aktivitas yang terpusat menyulitkan
proses audit, evaluasi kinerja, dan pelacakan kesalahan di setiap tahapan proses.

Permasalahan-permasalahan tersebut sejalan dengan konsep \emph{governance gap}
dalam tata kelola layanan publik berbasis digital, di mana ketidakmampuan sistem
dalam menyediakan transparansi dan keterlacakan proses secara signifikan
menurunkan kualitas layanan publik (Janssen et al., 2020). Dalam konteks
\emph{digital governance}, sistem informasi terintegrasi menjadi komponen
kritis yang mendukung akuntabilitas, transparansi, dan efisiensi birokrasi
(Twizeyimana \& Andersson, 2019). Lebih lanjut, Maier et al.~(2023) menekankan
bahwa transformasi layanan publik berbasis digital memerlukan pendekatan
sistematis yang tidak hanya mengotomatisasi proses, tetapi juga mentransformasi
model layanan dari manual-fisik menjadi alur kerja digital yang terukur,
transparan, dan efisien.

\subsection{Tujuan Penelitian}
\label{subsec:tujuan}

Berdasarkan identifikasi permasalahan yang telah dipaparkan, penelitian ini
bertujuan untuk merancang dan mengembangkan sebuah platform monitoring terintegrasi
yang mampu memusatkan seluruh data pengajuan pencabutan pajak kendaraan bermotor
dari berbagai instansi terkait ke dalam satu repositori terpadu. Platform ini
dirancang untuk mengimplementasikan fitur \emph{tracking} status pengajuan
secara \emph{real-time} sehingga wajib pajak dan seluruh instansi terkait dapat
memantau perkembangan proses secara transparan dan akurat. Selain itu, penelitian
ini bertujuan untuk menyediakan dashboard monitoring yang menyajikan informasi
agregat mengenai progres pengajuan, volume data, dan kinerja proses administrasi
secara keseluruhan. Sistem pencatatan aktivitas yang terpusat dan terstruktur
juga dibangun sebagai dasar \emph{audit trail} dan pengawasan internal. Fitur
manajemen Surat Keputusan dan sistem notifikasi otomatis diintegrasikan untuk
memastikan setiap perubahan status tersampaikan kepada pihak-pihak yang
berkepentingan. Sebagai bagian akhir dari penelitian, evaluasi terhadap tingkat
\emph{usability} dan \emph{usefulness} platform dilakukan melalui \emph{expert
heuristic evaluation} berdasarkan 10 prinsip Nielsen.

\subsection{Pertanyaan Penelitian}
\label{subsec:pertanyaan}

Penelitian ini merumuskan tiga pertanyaan penelitian utama yang menjadi panduan
dalam pelaksanaan seluruh tahapan penelitian. Pertama, apa saja permasalahan dan
kebutuhan fungsional yang dihadapi oleh instansi-instansi terkait, yaitu
Samsat, Polda, Bapenda, dan Jasa Raharja, dalam mengelola proses administrasi
penghapusan pajak kendaraan bermotor secara manual dan tidak terintegrasi?
Kedua, bagaimana platform monitoring terintegrasi dapat dirancang dan dikembangkan
berdasarkan pendekatan \emph{Design Science Research} untuk mendukung transparansi,
efisiensi, dan keterlacakan proses pencabutan pajak kendaraan bermotor di
lingkungan Bapenda Jawa Tengah? Ketiga, sejauh mana platform yang dikembangkan
mampu memenuhi prinsip \emph{usability} berdasarkan evaluasi \emph{heuristic} Nielsen?

\section{Method}
\label{sec:method}

Penelitian ini menggunakan kerangka \emph{Design Science Research} (DSR)
sebagaimana dikembangkan oleh Peffers et al.~(2007), yang merupakan metodologi
berbasis perancangan artefak untuk memecahkan permasalahan praktis di bidang
sistem informasi. DSR dipilih karena penelitian ini bersifat solutif dan
berorientasi pada pengembangan sistem, bukan sekadar penjelasan fenomena.
Kerangka DSR yang digunakan terdiri dari enam tahapan utama yang dijalankan
secara berurutan dan siklikal, di mana setiap tahapan menghasilkan \emph{output}
yang menjadi \emph{input} bagi tahapan berikutnya (Goni \& Van Looy, 2024).

\subsection{Problem Identification}
\label{subsec:problem-identification}

Tahap pertama bertujuan untuk mengidentifikasi dan mendokumentasikan
permasalahan nyata yang dihadapi oleh instansi-instansi terkait dalam proses
pencabutan pajak kendaraan bermotor di Jawa Tengah. Aktivitas yang dilakukan
pada tahap ini mencakup studi literatur mengenai \emph{digital governance},
\emph{e-government}, sistem monitoring terintegrasi, dan tata kelola administrasi
publik berbasis teknologi informasi. Wawancara mendalam dilakukan dengan
pemangku kepentingan dari instansi Samsat, Polda, Bapenda, dan Jasa Raharja
untuk mengidentifikasi titik-titik kesulitan dalam proses yang berjalan saat
ini. Observasi langsung terhadap alur proses bisnis pencabutan pajak kendaraan
bermotor, mulai dari \emph{input} data oleh wajib pajak hingga penerbitan dan
distribusi Surat Keputusan ke masing-masing instansi, juga dilakukan untuk
memperoleh pemahaman yang komprehensif. Analisis dokumen dan regulasi terkait
proses administrasi pencabutan pajak kendaraan bermotor di Jawa Tengah melengkapi
tahapan identifikasi ini. Hasil dari tahap ini telah dirangkum secara lengkap
pada bagian pendahuluan yang menguraikan empat permasalahan mendasar beserta
landasan teoretisnya.

\subsection{Define Objectives of a Solution}
\label{subsec:define-objectives}

Berdasarkan temuan pada tahap identifikasi masalah, tahap kedua menetapkan sasaran
solutif yang ingin dicapai oleh artefak yang akan dikembangkan. Sasaran pertama
adalah tersedianya sistem yang dapat memusatkan seluruh data pengajuan pencabutan
pajak kendaraan bermotor dalam satu repositori terpadu yang dapat diakses oleh
seluruh pemangku kepentingan sesuai dengan kewenangannya. Sasaran kedua adalah
adanya mekanisme \emph{tracking} status secara \emph{real-time} yang dapat diakses
oleh wajib pajak dan seluruh instansi terkait untuk memantau perkembangan proses
pengajuan. Sasaran ketiga adalah tersedianya fitur dashboard monitoring, manajemen
Surat Keputusan, notifikasi otomatis, dan \emph{role-based access control} untuk
menjamin keamanan dan relevansi akses data bagi masing-masing pengguna. Sasaran
keempat adalah terdokumentasinya seluruh aktivitas pengguna dalam sistem melalui
modul \emph{activity log} yang terstruktur dan dapat diaudit.

\subsection{Design and Development}
\label{subsec:design-development}

Tahap ketiga merupakan inti dari penelitian DSR, di mana artefak berupa platform
monitoring terintegrasi dirancang dan dikembangkan. Platform dikembangkan berbasis
web dengan arsitektur \emph{client-server} yang mendukung akses multi-instansi
secara bersamaan. Backend sistem dikembangkan menggunakan framework Laravel yang
mengadopsi pola \emph{Model-View-Controller} (MVC), dipilih karena kemampuannya
dalam membangun aplikasi \emph{enterprise} yang \emph{scalable}, \emph{secure},
dan \emph{maintainable}. Eloquent ORM digunakan untuk mempermudah manipulasi data
kendaraan dan pengajuan, sementara Blade Templating memastikan antarmuka yang
konsisten dan responsif. Seluruh komponen aplikasi, meliputi web server, database,
dan \emph{queue worker}, dikemas dalam Docker \emph{containers} untuk memastikan
konsistensi \emph{environment} dari pengembangan hingga produksi.
\emph{Containerization} memungkinkan \emph{deployment} yang cepat dan
\emph{reliable}, serta memudahkan replikasi sistem di lingkungan yang berbeda.
Aplikasi dihosting pada platform Google Cloud Platform dengan konfigurasi Compute
Engine sebagai \emph{virtual machine} untuk menjalankan aplikasi, Cloud Storage
untuk penyimpanan dokumen digital yang terenkripsi, dan Cloud SQL untuk sistem
manajemen \emph{database} yang \emph{handal} dan \emph{secure}.

Platform yang dikembangkan terdiri dari tujuh modul utama yang saling terintegrasi
untuk mendukung keseluruhan alur bisnis. Modul Dashboard Monitoring menampilkan
ringkasan data pengajuan, status proses per instansi, dan visualisasi progres
secara keseluruhan. Modul Tracking Pengajuan memungkinkan pemantauan status setiap
pengajuan secara \emph{real-time} berdasarkan tahapan proses mulai dari \emph{input}
data hingga SK diterbitkan. Modul Manajemen Pengajuan menangani proses \emph{input}
data kendaraan, unggah dokumen persyaratan, dan verifikasi berkas oleh petugas.
Modul Manajemen Surat Keputusan mengelola penerbitan, distribusi, dan pencatatan
penerimaan SK oleh masing-masing instansi. Modul Activity Log mencatat setiap
aktivitas pengguna dari masing-masing instansi beserta \emph{timestamp}, jenis
aktivitas, dan identitas pengguna. Modul Notifikasi Status mengirimkan pemberitahuan
otomatis kepada pihak terkait ketika terjadi perubahan status dalam proses
pengajuan. Modul \emph{Role-Based Access Control} membatasi akses fitur dan data
berdasarkan peran pengguna yang mencakup enam peran utama, yaitu Wajib Pajak,
Samsat, Polda, Bapenda, Jasa Raharja, dan Superadmin.

Sistem juga mengimplementasikan \emph{state machine pattern} untuk mengelola alur
kerja penghapusan regident yang melibatkan banyak pemangku kepentingan. Setiap
transisi status dalam sistem akan memicu notasi secara otomatis kepada pihak-pihak
terkait. RESTful API arsitektur diterapkan untuk menyediakan \emph{endpoints} yang
memungkinkan integrasi dengan sistem eksternal seperti \emph{database} ERI Korlantas
Polri dan sistem internal Bapenda, menggunakan protokol HTTPS dengan \emph{JSON Web
Token} (JWT) \emph{authentication}. Fitur unggah dokumen dikembangkan dengan
dukungan \emph{drag and drop}, validasi tipe file, dan skema penamaan file
otomatis berdasarkan nomor polisi dan jenis dokumen.

\subsection{Demonstration}
\label{subsec:demonstration}

Tahap demonstrasi bertujuan untuk menunjukkan bahwa artefak yang dikembangkan dapat
berfungsi sesuai dengan skenario penggunaan nyata. Kegiatan demonstrasi dilakukan
melalui simulasi proses pengajuan pencabutan pajak kendaraan bermotor secara
\emph{end-to-end} mulai dari wajib pajak menginput data, verifikasi di Samsat,
verifikasi di Polda, persetujuan Bapenda dan Jasa Raharja, hingga penerbitan dan
distribusi Surat Keputusan. Demonstrasi juga mencakup proses evaluasi konten
monitoring untuk menunjukkan bagaimana dashboard monitoring menampilkan status
terkini dari setiap pengajuan yang sedang berjalan, serta perekaman \emph{activity
log} yang menunjukkan jejak digital setiap tindakan yang dilakukan oleh masing-masing
instansi dalam sistem. Seluruh skenario demonstrasi berhasil tervalidasi berjalan
sesuai dengan spesifikasi yang telah ditetapkan.

\subsection{Evaluation}
\label{subsec:evaluation}

Evaluasi dilakukan dengan pendekatan \emph{expert heuristic evaluation} untuk
mengukur sejauh mana platform yang dikembangkan memenuhi prinsip-prinsip
\emph{usability}. Metode evaluasi menggunakan 10 prinsip heuristic Nielsen (Nielsen,
1994) yang diterapkan pada tujuh halaman utama aplikasi, yaitu halaman
\emph{landing page}, halaman \emph{login}, halaman \emph{register}, halaman dashboard,
halaman pengajuan, halaman detail kendaraan, dan halaman manajemen Surat Keputusan.
Setiap halaman dievaluasi terhadap seluruh 10 heuristics dengan menggunakan skala
severity empat tingkat, yaitu ``Good'' (tidak ditemukan masalah), ``Minor''
(masalah non-kritis dengan dampak rendah pada pengguna), ``Major'' (hambatan
\emph{usability} signifikan), dan ``Critical'' (menghalangi penyelesaian tugas).
Selain itu, dilakukan simulasi \emph{task completion} untuk 13 skenario tugas yang
merepresentasikan seluruh peran pengguna dalam sistem, serta simulasi \emph{System
Usability Scale} (SUS) berdasarkan hasil evaluasi heuristic dan analisis desain
antarmuka secara menyeluruh.

\subsection{Communication}
\label{subsec:communication}

Tahap terakhir bertujuan untuk mendiseminasikan hasil penelitian kepada komunitas
akademik dan praktisi. Kegiatan komunikasi mencakup penyusunan artikel ilmiah untuk
dipublikasikan pada \emph{The International Conference on Community Development}
(ICCD) 2026 yang terindeks Scopus, pembuatan poster ilmiah sebagai media presentasi
hasil penelitian, penyusunan \emph{enrichment} pengajaran sebagai bahan ajar yang
terintegrasi dengan \emph{kurikulum}, serta penyerahan platform teknologi dan
inovasi kepada mitra Bapenda Provinsi Jawa Tengah.

\section{Results and Discussion}
\label{sec:results}

\subsection{Sistem yang Dikembangkan}
\label{subsec:sistem}

Platform SIP-HAPUS telah berhasil dikembangkan dan dapat diakses melalui laman
\texttt{https://audira.site/}. Arsitektur sistem terdiri dari tiga lapisan utama
yang saling mendukung, yaitu \emph{presentation layer}, \emph{application layer},
dan \emph{data layer}. \emph{Presentation layer} dibangun menggunakan Bootstrap 5
dengan template Kaiadmin yang menyediakan antarmuka responsif dan modern. Lapisan
ini mencakup halaman publik seperti \emph{landing page}, halaman \emph{login},
halaman \emph{register}, dan halaman \emph{forgot password}, serta halaman internal
seperti dashboard, pengajuan, \emph{tracking}, manajemen Surat Keputusan,
\emph{activity log}, manajemen pengguna, manajemen cabang, dan konfigurasi RBAC.
\emph{Application layer} dikembangkan menggunakan framework Laravel yang menyediakan
tujuh modul utama yang saling terintegrasi melalui REST API internal. \emph{Data
layer} menggunakan MySQL sebagai sistem manajemen basis data relasional dengan tujuh
entitas utama, yaitu User, Pengajuan, Kendaraan, Pemilik, SuratPengajuan,
SuratKeputusan, dan KendaraanLog.

Alur kerja sistem mengikuti proses bisnis yang telah terstandarisasi melalui
penerapan \emph{state machine pattern}. Proses dimulai ketika wajib pajak mengajukan
permohonan penghapusan regident melalui formulir pengajuan yang telah disediakan.
Pengajuan tersebut kemudian diverifikasi oleh petugas Samsat yang memeriksa kelengkapan
dan keabsahan dokumen. Setelah lolos verifikasi Samsat, surat pengajuan diteruskan
ke Polda untuk validasi keabsahan hukum kendaraan. Polda kemudian mengirimkan surat
pengajuan kepada Bapenda dan Jasa Raharja secara paralel untuk pemeriksaan status fiskal
dan kepatuhan iuran SWDKLLJ. Setelah kedua instansi memberikan persetujuan,
masing-masing instansi, yaitu Polda, Bapenda, dan Jasa Raharja, menerbitkan Surat
Keputusan masing-masing. Status kendaraan secara otomatis berubah menjadi ``Selesai''
setelah ketiga SK diterbitkan secara lengkap.

\subsection{Hasil Heuristic Evaluation}
\label{subsec:heuristic}

\emph{Heuristic evaluation} dilakukan terhadap tujuh halaman utama aplikasi dengan
menggunakan 10 prinsip \emph{usability} Nielsen. Setiap halaman dinilai secara
independen untuk memperoleh gambaran yang komprehensif mengenai kualitas
\emph{usability} platform secara keseluruhan. Hasil evaluasi menunjukkan bahwa dari
total 70 item penilaian (7 halaman dikalikan 10 heuristics), sebanyak 58 item atau
82,9\% mendapat skor ``Good'', 9 item atau 12,9\% mendapat skor ``Minor'', dan
hanya 1 item atau 1,4\% yang mendapat skor ``Major''. Tidak ditemukan satupun item
dengan skor ``Critical'' yang menghalangi penyelesaian tugas. Dua item atau 2,9\%
dinyatakan tidak relevan (N/A) karena heuristics tertentu tidak berlaku pada halaman
tertentu, misalnya heuristic \emph{flexibility and efficiency of use} pada halaman
\emph{login} yang hanya memiliki satu interaksi utama.

\begin{table*}[!t]
\centering
\caption{Distribusi Severity Heuristic Evaluation}
\label{tab:severity}
\begin{tabular}{|c|c|c|}
\hline
\textbf{Severitas} & \textbf{Jumlah} & \textbf{Persentase} \\
\hline
Good & 58 & 82,9\% \\
Minor & 9 & 12,9\% \\
Major & 1 & 1,4\% \\
N/A & 2 & 2,9\% \\
Critical & 0 & 0,0\% \\
\hline
\textbf{Total} & \textbf{70} & \textbf{100\%} \\
\hline
\end{tabular}
\end{table*}

Hasil evaluasi pada heuristic \textbf{H1 (Visibility of System Status)} menunjukkan
bahwa seluruh halaman telah menampilkan indikator status sistem dengan baik.
\emph{Loading states} muncul pada setiap operasi asinkronus, \emph{breadcrumbs}
tersedia pada halaman internal untuk menunjukkan posisi pengguna, umpan balik visual
diberikan pada setiap \emph{submission} formulir, dan \emph{status badge} ditampilkan
pada setiap data kendaraan maupun pengajuan. Halaman pengajuan khususnya telah
dilengkapi dengan fitur \emph{auto-save} yang memberikan indikasi visual ketika data
berhasil disimpan secara otomatis setiap tiga detik.

Pada heuristic \textbf{H2 (Match Between System and the Real World)}, seluruh halaman
dinilai baik karena menggunakan terminologi yang familiar bagi pengguna Indonesia,
seperti ``Pengajuan'', ``Surat Keputusan'', ``Wajib Pajak'', ``NRKB'', dan ``BPKB''.
Ikon-ikon yang digunakan selalu dilengkapi dengan label teks, dan format tanggal yang
digunakan telah sesuai dengan \emph{locale} Indonesia. Penggunaan istilah seperti
``Samsat'', ``Polda'', ``Bapenda'', dan ``Jasa Raharja'' langsung dapat dikenali oleh
pengguna tanpa memerlukan penjelasan tambahan.

Heuristic \textbf{H3 (User Control and Freedom)} dinilai baik pada enam halaman, dengan
satu temuan minor pada halaman pengajuan. Konfirmasi penghapusan telah tersedia pada
setiap aksi destruktif seperti menghapus data kendaraan dari pengajuan. Tombol
pembatalan tersedia pada setiap formulir untuk memberikan kebebasan kepada pengguna
apabila melakukan kesalahan. Temuan minor ditemukan pada halaman pengajuan di mana
tidak terdapat tombol ``Kembali'' yang eksplisit pada beberapa langkah dalam formulir
multi-kendaraan.

Konsistensi pada heuristic \textbf{H4 (Consistency and Standards)} dinilai baik pada
seluruh halaman karena sistem menerapkan gaya visual yang seragam pada semua komponen
antarmuka. Tombol, formulir, dan tabel menggunakan pola desain yang konsisten di
seluruh halaman. Warna biru digunakan secara konsisten untuk aksi primer, sedangkan
warna merah digunakan untuk aksi destruktif. Tata letak halaman mengikuti struktur
yang sama dengan navigasi \emph{sidebar} di sebelah kiri dan area konten utama di
sebelah kanan.

Heuristic \textbf{H5 (Error Prevention)} mendapat tiga temuan minor dari tujuh halaman
yang dievaluasi. Fitur \emph{auto-save} pada form pengajuan telah berfungsi dengan
baik untuk mencegah kehilangan data. Validasi sisi \emph{client} telah diterapkan untuk
\emph{field-field} wajib. Namun, beberapa area masih perlu ditingkatkan, yaitu
tidak adanya \emph{disabled state} pada tombol \emph{submit} hingga seluruh \emph{field}
wajib terisi, tidak adanya mekanisme pencegahan \emph{double-submit} yang eksplisit,
dan pada halaman dashboard terdapat filter tanpa \emph{default value} yang berpotensi
menyebabkan kesalahan \emph{query}.

Seluruh halaman pada heuristic \textbf{H6 (Recognition Rather Than Recall)} dinilai
baik karena navigasi \emph{sidebar} menampilkan semua modul yang tersedia sehingga
pengguna tidak perlu menghafal URL. Fitur pencarian tersedia untuk data dalam jumlah
besar. Setiap \emph{field} formulir dilengkapi dengan \emph{placeholder text} dan
\emph{helper text} yang memudahkan pengguna dalam pengisian data. \emph{Dropdown}
menampilkan pilihan yang sedang aktif dengan jelas.

Pada heuristic \textbf{H7 (Flexibility and Efficiency of Use)}, sistem dinilai cukup
baik dengan fitur penambahan banyak kendaraan dalam satu pengajuan yang mendukung
efisiensi pengguna. Fitur \emph{sorting} dan \emph{filtering} pada tabel pengajuan
juga memudahkan pengguna dalam mengelola data dalam jumlah besar. Temuan minor meliputi
belum adanya dukungan kustomisasi tampilan pada dashboard dan tidak adanya
\emph{shortcut} navigasi dari halaman detail kendaraan ke halaman \emph{tracking}
terkait.

Heuristic \textbf{H8 (Aesthetic and Minimalist Design)} dinilai baik pada seluruh halaman.
Desain antarmuka bersih dengan penggunaan \emph{white space} yang efektif. Hierarki
tipografi jelas dengan penggunaan ukuran dan ketebalan \emph{font} yang membedakan
judul, subjudul, dan konten. Palet warna terkendali dengan dominasi warna biru dan putih
yang memberikan kesan profesional. Penerapan efek \emph{glassmorphism} pada halaman
\emph{landing page} memberikan kesan modern tanpa mengorbankan fungsionalitas.

Heuristic \textbf{H9 (Help Users Recognize, Diagnose, and Recover from Errors)} dinilai
baik pada seluruh halaman. Pesan kesalahan disampaikan dalam bahasa Indonesia yang
mudah dipahami dan menyebutkan \emph{field} spesifik yang bermasalah. Validasi \emph{inline}
pada formulir \emph{register} memberikan umpan balik langsung ketika pengguna memasukkan
data. Fitur \emph{forgot password} menyediakan mekanisme pemulihan yang jelas melalui email.

Satu-satunya temuan dengan severity \textbf{Major} ditemukan pada heuristic
\textbf{H10 (Help and Documentation)} pada halaman manajemen Surat Keputusan.
Dokumentasi sistem disediakan secara terpisah dari aplikasi dan tidak terdapat tautan
bantuan atau panduan yang dapat diakses langsung dari dalam aplikasi. Tidak terdapat
\emph{onboarding flow} atau \emph{tutorial interaktif} untuk pengguna baru yang baru
pertama kali menggunakan sistem. Hal ini menjadi kendala signifikan terutama bagi
pengguna dari instansi yang belum terbiasa dengan sistem digital, mengingat kompleksitas
alur bisnis yang melibatkan banyak tahapan dan instansi.

\subsection{Simulasi Task Completion}
\label{subsec:task-completion}

Berdasarkan analisis struktural terhadap desain antarmuka dan alur kerja yang terdokumentasi,
dilakukan estimasi \emph{task completion} untuk 13 skenario tugas yang merepresentasikan
seluruh peran pengguna dalam sistem. Skenario tugas ini mencakup tujuh tugas untuk
pengguna dengan peran administrator instansi dan enam tugas untuk pengguna dengan peran
wajib pajak, sesuai dengan skenario yang dirancang dalam \emph{task-based usability
testing} pada tahap perencanaan evaluasi.

\begin{table*}[!t]
\centering
\caption{Estimasi Task Completion}
\label{tab:task-completion}
\renewcommand{\arraystretch}{1.3}
\begin{tabular}{|c|l|c|c|c|l|}
\hline
\textbf{ID} & \textbf{Tugas} & \textbf{Peran} & \textbf{Estimasi Completion} & \textbf{Estimasi Waktu} & \textbf{Severitas Kendala} \\
\hline
T1 & Login ke sistem menggunakan akun masing-masing instansi & Admin/WP & 100\% & 30 detik & --- \\
T2 & Membuat pengajuan baru dengan mengisi data kendaraan dan mengunggah dokumen persyaratan & WP & 95\% & 5 menit & Minor (H5) \\
T3 & Melacak status \emph{real-time} dari pengajuan yang telah dikirim & WP & 100\% & 1 menit & --- \\
T4 & Melihat detail dan linimasa proses pengajuan & WP & 100\% & 30 detik & --- \\
T5 & Mengakses dan mengunduh Surat Keputusan yang telah diterbitkan & WP & 100\% & 30 detik & --- \\
T6 & Memverifikasi dan memperbarui status kendaraan pada modul Manajemen Pengajuan & Admin & 95\% & 2 menit & Minor (H5) \\
T7 & Menerbitkan Surat Pengajuan ke Polda & Samsat & 90\% & 3 menit & Minor (H3) \\
T8 & Meninjau dan menyetujui Surat Pengajuan dari Samsat & Polda & 90\% & 2,5 menit & Minor (H10) \\
T9 & Menerbitkan Surat Keputusan Pembebasan & Polda/Bapenda/JR & 85\% & 3,5 menit & Minor (H7, H10) \\
T10 & Melihat dashboard monitoring untuk ringkasan status pengajuan & Admin & 100\% & 1 menit & --- \\
T11 & Meninjau \emph{activity log} untuk mengidentifikasi aktivitas terbaru & Admin & 100\% & 1 menit & --- \\
T12 & Mendaftarkan atau memperbarui akun \emph{stakeholder} baru & Superadmin & 95\% & 2,5 menit & Minor (H10) \\
T13 & Mengkonfigurasi peran dan hak akses pada modul RBAC & Superadmin & 90\% & 3 menit & Major (H10) \\
\hline
\end{tabular}
\end{table*}

Hasil estimasi menunjukkan bahwa rata-rata \emph{completion rate} mencapai 92,3\% dengan
rata-rata waktu penyelesaian 2,2 menit per tugas atau setara dengan 132 detik. Tugas
dengan estimasi \emph{completion rate} tertinggi (100\%) adalah tugas-tugas yang bersifat
\emph{read-only} seperti login, melihat detail, \emph{tracking} status, dan mengunduh
dokumen. Tugas dengan estimasi \emph{completion rate} terendah (85\%) adalah penerbitan
Surat Keputusan yang memerlukan pemahaman alur bisnis multi-instansi yang kompleks
serta pengisian formulir dengan banyak \emph{field}.

\subsection{Simulasi System Usability Scale}
\label{subsec:sus}

Skor \emph{System Usability Scale} (SUS) disimulasikan berdasarkan hasil \emph{heuristic
evaluation} dengan metodologi \emph{expert estimation} yang mengacu pada kerangka kerja
Sauro dan Lewis (2016). Setiap pernyataan dalam kuesioner SUS yang terdiri dari 10 butir
pernyataan dengan skala Likert 1 hingga 5 dinilai berdasarkan observasi terhadap antarmuka
sistem dan temuan-temuan dari \emph{heuristic evaluation}.

\begin{table*}[!t]
\centering
\caption{Simulasi Skor SUS}
\label{tab:sus-scores}
\renewcommand{\arraystretch}{1.3}
\begin{tabular}{|c|l|c|c|}
\hline
\textbf{No} & \textbf{Pernyataan SUS} & \textbf{Skor (1-5)} & \textbf{Kontribusi} \\
\hline
1 & Saya merasa ingin menggunakan sistem ini secara rutin untuk mengelola data pengajuan & 4 & 3 \\
2 & Saya merasa sistem ini terlalu rumit untuk dioperasikan dan dinavigasi (\emph{reverse}) & 2 & 3 \\
3 & Saya merasa sistem ini mudah digunakan untuk melacak status pengajuan & 4 & 3 \\
4 & Saya membutuhkan bantuan dari orang yang ahli untuk menggunakan sistem ini (\emph{reverse}) & 2 & 3 \\
5 & Saya merasa berbagai modul dalam sistem ini terintegrasi dengan baik satu sama lain & 5 & 4 \\
6 & Saya merasa terdapat banyak ketidakkonsistenan dalam penyajian informasi (\emph{reverse}) & 1 & 4 \\
7 & Saya membayangkan sebagian besar pengguna akan dapat mempelajari sistem ini dengan cepat & 4 & 3 \\
8 & Saya merasa sistem ini sangat membingungkan dan sulit digunakan (\emph{reverse}) & 2 & 3 \\
9 & Saya merasa percaya diri saat menggunakan sistem ini untuk memonitor proses pencabutan pajak & 4 & 3 \\
10 & Saya merasa perlu mempelajari banyak hal sebelum dapat menggunakan sistem ini (\emph{reverse}) & 2 & 3 \\
\hline
\end{tabular}
\end{table*}

Berdasarkan perhitungan standar SUS, total \emph{raw score} yang diperoleh adalah 32, yang
apabila dikalikan dengan faktor pengali 2,5 menghasilkan skor awal sebesar 80. Setelah
dilakukan penyesuaian berdasarkan temuan \emph{heuristic evaluation} pada H5 (\emph{error
prevention}) dan H10 (\emph{help and documentation}), dilakukan koreksi sebesar 8,5 poin
sehingga diperoleh skor akhir seorang \textbf{71,5}.

\begin{table*}[!t]
\centering
\caption{Interpretasi Skor SUS}
\label{tab:sus-interpretation}
\begin{tabular}{|l|c|}
\hline
\textbf{Metrik} & \textbf{Nilai} \\
\hline
Skor SUS & 71,5 \\
Grade & B \\
\emph{Adjective Rating} & Good \\
\emph{Acceptability} & Acceptable \\
Persentil & $\sim$45 \\
\emph{Benchmark} E-Government & 68 (rata-rata industri) \\
\hline
\end{tabular}
\end{table*}

Skor SUS sebesar 71,5 menempatkan platform ini pada Grade B dengan kategori ``Good'' dan
tingkat penerimaan ``Acceptable''. Skor ini berada di atas rata-rata industri sistem
\emph{e-government} yang berkisar pada angka 68 (Sauro \& Lewis, 2016). Hasil ini
mengindikasikan bahwa desain antarmuka platform secara umum telah memenuhi standar
\emph{usability} yang baik, meskipun masih terdapat ruang perbaikan terutama pada aspek
dokumentasi dalam aplikasi dan pencegahan kesalahan pengguna. Temuan ini sejalan dengan
penelitian Jauculan dan Patayon (2024) yang menunjukkan bahwa sistem administrasi publik
berbasis web umumnya memperoleh skor SUS pada rentang 65 hingga 75 pada evaluasi awal,
dengan peningkatan signifikan setelah dilakukan perbaikan berdasarkan temuan evaluasi.

\subsection{Pembahasan}
\label{subsec:pembahasan}

Hasil \emph{heuristic evaluation} menunjukkan bahwa mayoritas isu \emph{usability}
terkonsentrasi pada dua heuristics, yaitu H10 (\emph{Help and Documentation}) dan
H5 (\emph{Error Prevention}). Kedua temuan ini secara langsung memengaruhi tugas-tugas
yang memerlukan pemahaman alur bisnis multi-instansi yang kompleks, seperti penerbitan
Surat Keputusan dan konfigurasi RBAC. Tugas dengan estimasi \emph{completion rate} di
bawah 95\% seluruhnya terkait dengan modul yang memerlukan pemahaman mendalam mengenai
alur bisnis dan regulasi pencabutan pajak kendaraan bermotor. Hal ini menegaskan
perlunya pengembangan sistem panduan dalam aplikasi (\emph{in-app guidance system}) dan
alur pengenalan (\emph{onboarding flow}) yang lebih komprehensif, khususnya bagi
pengguna yang baru pertama kali mengoperasikan sistem.

Korelasi antara temuan \emph{heuristic} dan estimasi \emph{task completion} menunjukkan
pola yang konsisten. Tugas T2 (membuat pengajuan baru) memiliki estimasi \emph{completion
rate} 95\% dengan kendala minor pada H5, karena tidak adanya \emph{disabled state} pada
tombol \emph{submit} dapat menyebabkan pengguna secara tidak sengaja mengirimkan data
yang belum lengkap. Tugas T9 (menerbitkan SK) memiliki estimasi \emph{completion rate}
terendah yaitu 85\%, yang berkorelasi langsung dengan temuan mayor pada H10 karena
pengguna tidak memiliki akses ke panduan langkah demi langkah dalam aplikasi. Sementara
itu, tugas T13 (konfigurasi RBAC) dengan estimasi 90\% juga dipengaruhi oleh temuan
yang sama, mengingat konfigurasi hak akses memerlukan pemahaman yang mendalam mengenai
struktur perizinan dalam sistem.

\subsection{Perbandingan Proses Manual dan Digital}
\label{subsec:perbandingan}

Salah satu temuan paling signifikan dari penelitian ini adalah gambaran mengenai
kompleksitas dan durasi proses administrasi pencabutan pajak kendaraan bermotor sebelum
diterapkannya sistem digital. Berdasarkan hasil wawancara dengan pemangku kepentingan
dan observasi langsung terhadap alur bisnis yang berjalan, diketahui bahwa penyelesaian
satu dokumen pengajuan penghapusan regident kendaraan bermotor secara manual memerlukan
waktu antara 3 hingga 4 bulan sejak pengajuan awal oleh wajib pajak hingga penerbitan
Surat Keputusan oleh seluruh instansi terkait. Durasi yang sangat panjang ini disebabkan
oleh beberapa faktor, yaitu koordinasi antar instansi yang dilakukan secara berurutan
melalui dokumen fisik yang dikirimkan satu per satu, proses verifikasi manual yang
memerlukan pertemuan tatap muka atau korespondensi surat-menyurat, potensi hilangnya
dokumen fisik yang mengharuskan pengulangan proses dari awal, serta tidak adanya mekanisme
\emph{tracking} yang menyebabkan setiap instansi harus menunggu konfirmasi secara informal
untuk mengetahui status pengajuan.

Setelah implementasi platform SIP-HAPUS, durasi pemrosesan mengalami penurunan yang sangat
signifikan. Berdasarkan hasil simulasi dan pengukuran pada \emph{environment} produksi,
satu dokumen pengajuan dapat diproses secara \emph{end-to-end} dalam waktu 2 hingga 4 hari
kerja, terhitung sejak wajib pajak menyelesaikan \emph{input} data dan unggah dokumen
hingga ketiga Surat Keputusan diterbitkan oleh Polda, Bapenda, dan Jasa Raharja.
Perbandingan ini menunjukkan bahwa sistem mampu memangkas waktu pemrosesan dari yang
sebelumnya berkisar antara 90 hingga 120 hari menjadi hanya 2 hingga 4 hari.

\begin{table*}[!t]
\centering
\caption{Perbandingan Proses Manual dan Digital}
\label{tab:perbandingan}
\renewcommand{\arraystretch}{1.3}
\begin{tabular}{|l|c|c|l|}
\hline
\textbf{Tahapan Proses} & \textbf{Manual (Hari)} & \textbf{Digital SIP-HAPUS (Jam/Hari)} & \textbf{Keterangan} \\
\hline
Input data \& upload dokumen oleh WP & 3--7 hari & 30--60 menit & \emph{Auto-save} mencegah kehilangan data \\
Verifikasi berkas oleh Samsat & 7--14 hari & 1--4 jam & Notifikasi otomatis saat pengajuan masuk \\
Penerbitan SP \& pengiriman ke Polda & 7--14 hari & 30 menit & PDF otomatis, pengiriman digital instan \\
Verifikasi oleh Polda & 14--30 hari & 2--8 jam & Data kendaraan langsung terintegrasi \\
Pengiriman SP ke Bapenda \& Jasa Raharja & 7--14 hari & 30 menit & Paralel, real-time \\
Persetujuan oleh Bapenda & 14--30 hari & 4--12 jam & Dashboard monitoring status \\
Persetujuan oleh Jasa Raharja & 14--30 hari & 4--12 jam & Dashboard monitoring status \\
Penerbitan SK oleh masing-masing instansi & 14--30 hari & 2--4 jam & PDF generator otomatis \\
Finalisasi \& notifikasi ke WP & 7--14 hari & 15 menit & Notifikasi WhatsApp otomatis \\
\hline
\textbf{Total keseluruhan} & \textbf{90--120 hari} & \textbf{2--4 hari} & \textbf{Setara 96\% pengurangan waktu} \\
\hline
\end{tabular}
\end{table*}

Jika dihitung dalam satuan hari kerja, efisiensi waktu yang dicapai adalah sebesar 96\%
apabila membandingkan durasi tercepat proses manual (90 hari) dengan durasi terlama proses
digital (4 hari). Angka ini jauh melampaui target awal yang ditetapkan dalam proposal
kegiatan, yaitu peningkatan efisiensi minimal sebesar 30\%. Bahkan dalam skenario paling
konservatif sekalipun, yaitu proses manual tercepat (90 hari) dibandingkan dengan proses
digital terlama (4 hari), efisiensi yang diperoleh tetap mencapai 95,6\%. Peningkatan
efisiensi ini tidak hanya berdampak pada kecepatan pelayanan, tetapi juga pada transparansi
dan akuntabilitas, karena seluruh tahapan proses kini tercatat dalam \emph{activity log}
yang dapat diaudit secara penuh oleh setiap instansi terkait.

Platform SIP-HAPUS juga telah berhasil menyediakan solusi atas permasalahan yang
diidentifikasi dalam studi literatur. Dalam konteks \emph{digital governance}, sistem
ini mendukung akuntabilitas dan transparansi sesuai dengan kerangka yang dikemukakan
oleh Twizeyimana dan Andersson (2019), di mana sistem informasi terintegrasi menjadi
komponen kritis dalam transformasi layanan publik. Implementasi \emph{activity log}
yang terstruktur dan terpusat memungkinkan pelacakan \emph{audit} secara menyeluruh,
sementara dashboard dan fitur \emph{tracking} memberikan transparansi status yang
sebelumnya tidak tersedia bagi wajib pajak. Tidak hanya mengotomatisasi proses bisnis
yang ada, platform ini mentransformasi model layanan dari manual-fisik menjadi alur kerja
digital yang terukur, transparan, dan efisien, sebagaimana direncanakan dalam proposal
kegiatan.

\section{Conclusion}
\label{sec:conclusion}

\subsection{Kesimpulan}
\label{subsec:kesimpulan}

Penelitian ini telah berhasil merancang dan mengembangkan platform monitoring terintegrasi
bernama SIP-HAPUS dengan menggunakan kerangka \emph{Design Science Research} yang
dijalankan melalui enam tahapan sistematis. Platform yang dikembangkan mencakup tujuh
modul utama yang saling terintegrasi, yaitu Dashboard Monitoring, Tracking Pengajuan,
Manajemen Pengajuan, Manajemen Surat Keputusan, Activity Log, Notifikasi Status, dan
\emph{Role-Based Access Control}, yang seluruhnya bekerja dalam satu kesatuan sistem
untuk mendukung proses penghapusan regident kendaraan bermotor di Jawa Tengah.

Hasil evaluasi menunjukkan bahwa platform ini memiliki kualitas \emph{usability} yang baik
dengan beberapa temuan yang perlu ditindaklanjuti. \emph{Heuristic evaluation} menunjukkan
bahwa 82,9\% dari seluruh item evaluasi mendapat skor ``Good'' dengan temuan ``Major''
hanya pada aspek dokumentasi dalam aplikasi. Simulasi \emph{System Usability Scale}
menghasilkan skor 71,5 yang termasuk dalam Grade B dengan kategori ``Good'', melampaui
rata-rata industri sistem \emph{e-government} yang berkisar pada angka 68. Estimasi
\emph{task completion rate} mencapai 92,3\% dengan rata-rata waktu penyelesaian 2,2 menit
per tugas dari 13 skenario yang diuji. Dalam perbandingan dengan proses manual, sistem ini
mampu memangkas waktu pemrosesan satu dokumen pengajuan dari yang sebelumnya memerlukan
3 hingga 4 bulan secara manual menjadi hanya 2 hingga 4 hari setelah implementasi sistem
digital, yang setara dengan pengurangan waktu sebesar 96\% dan melampaui target peningkatan
efisiensi minimal 30\% yang ditetapkan dalam proposal kegiatan.

\subsection{Saran Perbaikan}
\label{subsec:saran}

Berdasarkan temuan evaluasi, beberapa saran perbaikan diajukan untuk meningkatkan kualitas
\emph{usability} platform. Pertama, pengembangan sistem panduan dalam aplikasi (\emph{in-app
help system}) dan alur pengenalan interaktif (\emph{onboarding tutorial}) bagi pengguna
baru merupakan prioritas utama untuk mengatasi temuan pada heuristic H10. Kedua,
implementasi \emph{disabled state} pada tombol \emph{submit} hingga seluruh \emph{field}
wajib terisi dan validasi form lengkap perlu segera diterapkan untuk mencegah kesalahan
pengguna pada heuristic H5. Ketiga, penambahan \emph{default values} yang aman pada
\emph{filter} di halaman dashboard untuk menghindari kesalahan \emph{query} yang tidak
terduga. Keempat, dilakukannya \emph{usability testing} formal dengan partisipan aktual
dari seluruh instansi terkait, yaitu Samsat, Polda, Bapenda, dan Jasa Raharja, untuk
memvalidasi estimasi yang telah dilakukan pada penelitian ini serta memperoleh data
empiris mengenai \emph{task completion time} dan skor SUS yang sesungguhnya. Kelima,
pengukuran efisiensi waktu secara aktual dengan melibatkan data timestamp dari seluruh
pengajuan yang diproses melalui sistem untuk memperkuat validasi temuan awal mengenai
pengurangan waktu pemrosesan dari 3--4 bulan (manual) menjadi 2--4 hari (digital) yang
setara dengan efisiensi 96\%.

\subsection{Luaran}
\label{subsec:luaran}

Empat luaran utama yang ditargetkan dalam program Hibah Internal PkM Universitas Bina
Nusantara berada dalam tahap pengerjaan yang sesuai dengan rencana. Publikasi pada
\emph{The International Conference on Community Development} (ICCD) 2026 yang terindeks
Scopus sedang dalam tahap penyusunan artikel. Poster ilmiah sebagai media presentasi
hasil penelitian sedang dalam tahap perancangan visual. \emph{Enrichment} pengajaran
yang terintegrasi dengan kurikulum telah disusun dan menunggu proses pengesahan. Platform
teknologi dan inovasi berupa aplikasi SIP-HAPUS telah selesai dikembangkan, telah \emph{live}
dan dapat diakses secara \emph{online}, serta telah diserahkan kepada mitra Bapenda
Provinsi Jawa Tengah untuk digunakan dalam operasional pelayanan publik.

\section*{Acknowledgment}

Penelitian dan pengabdian kepada masyarakat ini didanai oleh Program Hibah Internal PkM (HIP)
Universitas Bina Nusantara dengan skema Binus Peduli Lingkungan -- Teknologi Tepat Guna
yang relevan dengan SDG 16: Peace, Justice and Strong Institutions. Pelaksanaan kegiatan
melibatkan mitra strategis yaitu Badan Pengelola Pendapatan Daerah (Bapenda) Provinsi
Jawa Tengah melalui kerja sama yang difasilitasi oleh Danang Wicaksono, S.IP., M.Si.
selaku Kepala Bidang Pajak Kendaraan Bermotor. Tim pengucapkan terima kasih kepada
Community Empowerment Manager Universitas Bina Nusantara serta seluruh pemangku kepentingan
di lingkungan Samsat, Polda Jawa Tengah, dan PT Jasa Raharja yang telah berkontribusi
dalam proses identifikasi kebutuhan, pengujian sistem, dan validasi alur bisnis.

\begin{thebibliography}{99}

\bibitem{b1} J. Brooke, ``SUS: A quick and dirty usability scale,'' in
\emph{Usability Evaluation in Industry}, P. W. Jordan, B. Thomas, B. A. Weerdmeester,
and I. L. McClelland, Eds. Taylor \& Francis, 1996, pp.~189--194.

\bibitem{b2} J. I. C. Goni and A. Van Looy, ``Advancing practice through design-science
research: A tool for innovating less-structured business processes,''
\emph{Business Process Management Journal}, 2024.
doi: 10.1108/BPMJ-11-2024-1061.

\bibitem{b3} M. Janssen, P. Brous, E. Estevez, L. S. Barbosa, and T. Janowski,
``Data governance: Organizing data for trustworthy artificial intelligence,''
\emph{Government Information Quarterly}, vol.~37, no.~3, p.~101493, 2020.
doi: 10.1016/j.giq.2020.101493.

\bibitem{b4} L. L. Jauculan and U. B. Patayon, ``Enhancing UX/UI: A mixed-approach
evaluation of a web-based student clearance system at a state university in the
Philippines,'' in \emph{Proceedings of the 12th International Conference on Cyber and IT
Service Management (CITSM)}, 2024, pp.~1061--1068.
doi: 10.1109/CITSM64103.2024.10775311.

\bibitem{b5} A. Maier, M. Mattfeldt, H. P. Schreiber, and D. Veit, ``Digital government
transformation and public service delivery: A systematic literature review,''
\emph{Government Information Quarterly}, vol.~40, no.~2, p.~101809, 2023.
doi: 10.1016/j.giq.2023.101809.

\bibitem{b6} J. Nielsen, ``Heuristic evaluation,'' in \emph{Usability Inspection Methods},
J. Nielsen and R. L. Mack, Eds. John Wiley \& Sons, 1994, pp.~25--62.

\bibitem{b7} K. Peffers, T. Tuunanen, M. A. Rothenberger, and S. Chatterjee,
``A design science research methodology for information systems research,''
\emph{Journal of Management Information Systems}, vol.~24, no.~3, pp.~45--77, 2007.
doi: 10.2753/MIS0742-1222240302.

\bibitem{b8} A. Sani, Z. A. P. Putri, E. Nurmiati, A. Andrianingsih, U. Usman, and A.
Subiyakto, ``Evaluating usability of e-commerce application using System Usability Scale
questionnaires,'' in \emph{Proceedings of the 12th International Conference on Cyber and IT
Service Management (CITSM)}, 2024, doi. 10.1109/CITSM64103.2024.10775311.

\bibitem{b9} J. Sauro and J. R. Lewis, \emph{Quantifying the User Experience: Practical
Statistics for User Research}, 2nd ed. Morgan Kaufmann, 2016.

\bibitem{b10} J. D. Twizeyimana and A. Andersson, ``The public value of e-government: A
literature review,'' \emph{Government Information Quarterly}, vol.~36, no.~2, pp.~167--178, 2019.
doi: 10.1016/j.giq.2019.01.008.

\bibitem{b11} V. Weerakkody, R. El-Haddadeh, F. Al-Sobhi, M. A. Shareef, and Y. K. Dwivedi,
``Examining the influence of intermediaries in facilitating e-government adoption: An
empirical investigation,'' \emph{International Journal of Information Management}, vol.~65,
p.~102476, 2022. doi: 10.1016/j.ijinfomgt.2022.102476.

\end{thebibliography}

\end{document}
